% Generated draft: requires mathematical and source review.
\documentclass[11pt]{article}
\usepackage[margin=25mm]{geometry}
\usepackage{fontspec}
\setmainfont{Latin Modern Roman}
\usepackage{amsmath,amssymb,mathtools,unicode-math}
\setmathfont{Latin Modern Math}
\usepackage{enumitem,needspace,xurl,hyperref}
\setlength{\parindent}{0pt}
\setlength{\parskip}{5pt}
\setlength{\emergencystretch}{3em}
\newcommand{\N}{\mathbb{N}}
\newcommand{\Z}{\mathbb{Z}}
\newcommand{\Q}{\mathbb{Q}}
\newcommand{\R}{\mathbb{R}}
\newcommand{\C}{\mathbb{C}}
\newcommand{\F}{\mathbb{F}}
\newcommand{\A}{\mathbb{A}}
\newcommand{\PP}{\mathbb{P}}
\newcommand{\E}{\mathbb{E}}
\newcommand{\eps}{\varepsilon}
\newcommand{\dd}{\,\mathrm{d}}
\begin{document}
\section*{UnsolvedMath ID: 30001285}
\textbf{Source problem code:} catalogue entry
\subsection{Definitions and mathematical statement}
Let \(k\) denote the unspecified base field and let \(A\) be a central simple \(k\)-algebra: a finite-dimensional associative unital \(k\)-algebra with center \(k\) and no nonzero proper two-sided ideals. Its reduced norm defines the algebraic group
\[
\operatorname{SL}_1(A)=\ker\bigl(\operatorname{Nrd}_{A/k}:\operatorname{GL}_1(A)\longrightarrow\mathbb G_m\bigr).
\]
Generalized Severi--Brauer varieties associated with \(A\) parametrize ideals of prescribed reduced dimension; the supplied record does not specify the dimension parameters or the ideal convention involved.

The source asks for a comparison of three motivic-cohomology maps associated with \(A\), arising respectively from the \'{e}tale Bloch--Lichtenbaum spectral sequence, generalized Severi--Brauer varieties, and \(\operatorname{SL}_1(A)\). Under which hypotheses, and after which identifications of their sources and targets, do these constructions give the same map?

More explicitly, if the intended constructions and identifications produce homomorphisms
\[
f_{\mathrm{et}},\ f_{\mathrm{SB}},\ f_{\mathrm{SL}_1}:M\longrightarrow N,
\]
the requested comparison is whether, and under what conditions,
\[
f_{\mathrm{et}}=f_{\mathrm{SB}}=f_{\mathrm{SL}_1}.
\]
Here \(M\), \(N\), and the three homomorphisms are schematic placeholders, not objects specified by the supplied record. The record omits their definitions, cohomological degrees, coefficients, and comparison identifications, as well as any restrictions on \(k\) or \(A\). Consequently, it does not determine a precise equality of defined maps or a fully quantified conjecture. These missing data require source clarification and human mathematical review.
\subsection{Short English statement}
For a central simple algebra, compare the motivic-cohomology maps obtained from the \'{e}tale Bloch--Lichtenbaum spectral sequence, generalized Severi--Brauer varieties, and \(\operatorname{SL}_1(A)\). The supplied question asks when they agree, but does not define the maps or their common source and target, so a precise equality remains unresolved.
\subsection{Sources}
\begin{itemize}
\item Local UnsolvedMath registry, permanent ID \(30001285\), source problem code \texttt{catalogue entry}, title: \emph{Comparison Maps in Motivic Cohomology of Central Simple Algebras}. The supplied record belongs to the Oberwolfach Reports Open Problems collection and gives proposed year \(2009\); no individual proposer is identified.
\item \emph{Algebraic \(K\)-Theory and Motivic Cohomology} (2009), Oberwolfach Reports, pp.~1731--1774, DOI \texttt{10.4171/owr/2009/31}. Supplied source report and literature-source link, labeled ``OWR report'': \url{https://doi.org/10.4171/owr/2009/31}. The background describes these reports as documenting mathematical workshops and their open-problem discussions.
\item Literature assessment dated \texttt{2026-08-21}: the supplied curation labels the problem open, reports that the report explicitly calls comparison of the maps an open problem, and states that no general comparison theorem was found. This is supplied dated curation, not an independently verified literature assessment; time-sensitive status and cited claims require independent verification.
\item The registry labels the statement \texttt{corrected\_verified} and states that its mathematical content was checked against the official report, with notation, context, and question form normalized without changing the source-level claim. This verification claim is supplied provenance and has not been independently checked for this draft.
\end{itemize}
\end{document}
